\begin{abstract}
Radio-frequency transmitter recognition can depend on the receiver used to acquire a signal, complicating transfer to unseen hardware. We compare source-only classification, test-time classifier adjustment, and learned receiver-context conditioning on a frozen six-transmitter WiSig task with leave-one-receiver-out evaluation across 32 physical receivers and five seeds. Target information is bounded by a label-free, deterministic bank of 128 receiver-specific packets disjoint from the query pool. Frozen results give receiver-equal macro-F1 of \val{P0}{macro_f1}{raw} for the independent baseline, \val{T3A}{macro_f1}{raw} for T3A, and \val{P2}{macro_f1}{raw} for attentive context conditioning. The historical T3A-minus-baseline difference is positive for \dval{T3A_MINUS_P0}{positive} receivers, with receiver-bootstrap interval [\dval{T3A_MINUS_P0}{ci95_lower}, \dval{T3A_MINUS_P0}{ci95_upper}]. A separately frozen, explicitly post-hoc reviewer-completeness addendum evaluates GroupNorm-compatible sharpness-aware adaptation, source-aligned embedding standardization, and full-network labeled adaptation. Embedding standardization reaches \val{EMB_STD}{macro_f1}{raw}; the sharpness-aware application remains near the independent baseline. Full-network labeled adaptation reaches \val{SUP_FT_FULL_128}{macro_f1}{raw}, invalidating interpretation of the earlier head-only oracle as a strong ceiling. Richer probability analysis identifies average underconfidence in raw T3A despite improved proper scores. Source-only temperature fitting improves probability quality without changing decisions. The evidence supports a bounded single-dataset receiver-adaptation benchmark, not external generality or physically acquired adaptation episodes.
\end{abstract}
